\section*{第 \ref{ch:b1:alcohol-activation} 章——醇：活化 OH 基团}

\begin{solution}{exo:b1:alcohol-activation:1}
\ce{2CH3CH2OH + 2Na -> 2CH3CH2O- + 2Na+ + H2};
\ce{CH3CH2OH + NaH -> CH3CH2O- + Na+ + H2};
\ce{CH3CH2OH + OH- <=> CH3CH2O- + H2O}, $K = 10^{14.0 - 15.9} = 10^{-1.9}$:
仅部分进行。
\end{solution}

\begin{solution}{exo:b1:alcohol-activation:2}
1-甲氧基丙烷；甲氧基乙烷；乙氧基苯（苯氧负离子，
碱性弱于醇盐，仍是良好的亲核试剂）。
\end{solution}

\begin{solution}{exo:b1:alcohol-activation:3}
在吡啶中进行对甲苯磺酰化，得到 \ce{CH3CH2CH2OTs}；随后
\[ \ce{CH3CH2CH2OTs + I- -> CH3CH2CH2I + TsO-}. \] 若直接用碘离子，则须
逐出 \ce{OH-}，即不反应。对甲苯磺酸酯具有很好的离去基团，
且条件既非酸性，亦非强碱性。
\end{solution}

\begin{solution}{exo:b1:alcohol-activation:4}
丁-2-醇：丁-1-烯、\cip{Z}-丁-2-烯和
\cip{E}-丁-2-烯，以后者为主。2-甲基戊-2-醇：2-甲基戊-2-烯（主产物，
三取代）和 2-甲基戊-1-烯。
\end{solution}

\begin{solution}{exo:b1:alcohol-activation:5}
\ce{(CH3)2CHCH2-O-CH2CH3}：2-甲基丙-1-醇钠与溴乙烷反应，
或乙醇钠与 1-溴-2-甲基丙烷反应。两种卤代烃均为伯卤代烃，
但 1-溴-2-甲基丙烷在反应碳旁带有支链，
这会减慢 SN2 反应，并使强碱发生消除生成 2-甲基丙烯。
前一方案更佳。
\end{solution}

\begin{solution}{exo:b1:alcohol-activation:6}
对甲苯磺酸酯保持 \cip{R} 构型；与乙醇钠发生 SN2 反应时构型翻转，
得 \cip{S}-2-乙氧基辛烷（OEt 与 OTs 同样优先）；用硫酸并加热时，
发生消除生成辛烯（主要为 \cip{E}-辛-2-烯），经由
仲碳正离子，立体构型随之丧失。
\end{solution}

\begin{solution}{exo:b1:alcohol-activation:7}
羟基质子化；失水生成叔碳正离子（慢步）；
被 \ce{Cl-} 捕获，生成 2-氯-2-甲基丙烷。该步很快，原因是叔
碳正离子稳定，且氯化物在酸中不溶而分离。
\end{solution}

\begin{solution}{exo:b1:alcohol-activation:8}
\ce{CH3CH2OH + H+ -> CH3CH2OH2+}；第二个乙醇分子从背面进攻该碳，
水离去（SN2），生成 \ce{(CH3CH2)2OH+}，它失去质子：
乙氧基乙烷。竞争反应中，碱（乙醇、\ce{HSO4-}）夺取 $\beta$
位质子（\ce{CH3CH2OH2+}），水离去时发生 E2 消除，生成乙烯。
\end{solution}

\begin{solution}{exo:b1:alcohol-activation:9}
2-甲基丙-2-醇（叔醇），生成难溶的 2-氯-2-甲基丙烷。
\end{solution}

\begin{solution}{exo:b1:alcohol-activation:10}
质子化并失水，生成仲碳正离子，它紧邻
一个季碳。该季碳上的一个甲基连同其成键电子对
迁移至带正电的碳上，正电荷随之位于叔碳上，
更稳定。邻位碳失去质子，
生成四取代的 2,3-二甲基丁-2-烯，即最稳定的烯烃。
\end{solution}

\begin{solution}{exo:b1:alcohol-activation:11}
由 \cip{R}-丁-2-醇出发：先得对甲苯磺酸酯（\cip{R}），再与甲醇钠反应（SN2，
构型翻转），得 \cip{S}-2-甲氧基丁烷。由 \cip{S}-丁-2-醇出发，先制成醇盐
（NaH），再加碘甲烷：立体中心碳的 C--O 键未被破坏，
故仍得 \cip{S}-2-甲氧基丁烷。
\end{solution}

\begin{solution}{exo:b1:alcohol-activation:12}
$Q = [\text{醚}][\ce{H2O}]/[\ce{EtOH}]^2$。随生成随移除水，
$[\ce{H2O}]$ 和 $Q$ 保持较小，$Q < K$，反应继续生成
醚，直至醇耗尽。
\end{solution}

\begin{solution}{pb:b1:alcohol-activation:1}
\textbf{1.} 叔丁醇钠与卤甲烷反应；甲醇钠与
叔丁基卤反应。
\textbf{2.} 第二种：甲醇盐是强碱，会使
叔卤代烷发生消除：\ce{(CH3)3CBr + CH3O- -> (CH3)2C=CH2 + CH3OH + Br-}。
\textbf{3.} \ce{(CH3)3CO- + CH3I -> (CH3)3COCH3 + I-}：醇盐的氧原子
从背面进攻甲基碳，碘离子同时离去（SN2）。
\textbf{4.} 用钠（或氢化钠）反应：\ce{2(CH3)3COH + 2Na ->
2(CH3)3CO- + 2Na+ + H2}。氢氧根碱性过弱：$K = 10^{14.0 -
19.2} = 10^{-5.2}$.
\textbf{5.} 水会使醇盐质子化，并进攻卤代烷。
\textbf{6.} 否：不存在立体中心。
\textbf{7.} \cip{R}-丁-2-醇 + TsCl（吡啶）$\to$ \cip{R}-丁-2-基
对甲苯磺酸酯：构型保持。
\textbf{8.} \cip{S}-2-甲氧基丁烷，经 SN2 反应构型翻转。
\textbf{9.} \cip{S}-丁-2-醇。
\textbf{10.} \cip{R}-2-甲氧基丁烷：立体碳的 C--O 键
保持不变，只形成 O--C（甲基）键。
\textbf{11.} E2（甲醇盐为强碱，且碳为仲碳），
生成丁烯；消除程度受碱较小、离去基团极好，
以及温度温和所限。
\textbf{12.} 质子化的醇电离为平面的仲碳
正离子，两面均被甲醇捕获：得外消旋醚，
并伴生丁烯。
\textbf{13.} 羟基质子化，失水（慢步）生成叔
碳正离子，再失去一个 $\beta$ 质子。
\textbf{14.} 2-甲基丁-2-烯（主产物，三取代）和
2-甲基丁-1-烯。
\textbf{15.} 叔碳正离子受阻，失去质子
远比另一醇分子对其进攻更容易。
\textbf{16.} 移除产物：$Q$ 保持较小，反应得以继续，
且烯烃不致与酸长时间接触。
\textbf{17.} 生成碳正离子的第一能垒（失水，决速）
较高，随后生成烯烃的能垒较小。
\textbf{18.} 否：伯碳正离子过于不稳定；伯醇
在质子化的醇上经 E2 消除，需较高温度。
\textbf{19.} $10.0/74.0 = \qty{0.135}{mol}$.
\textbf{20.} $1.10 \times 0.135 \times 141.9 = \qty{21.1}{g}$.
\textbf{21.} $0.135 \times 88.0 = \qty{11.9}{g}$.
\textbf{22.} $0.75 \times 11.9 =$ \textbf{\qty{8.9}{g} 的 MTBE}。
\end{solution}
